\section{A Complete Plonk Example over \texorpdfstring{$\F_{97}$}{F97}}

This appendix gives a complete Plonk-style proof sketch at the arithmetization and IOP-gadget level. It includes the ProductCheck inside the wiring/permutation argument. It is not yet a production Plonk transcript with all batching, blinding, and Fiat-Shamir challenge ordering spelled out.

\subsection{Field and subgroup}

Work over the prime field
\[
    \F_{97}.
\]
We use a subgroup of size $12$ because the trace occupies $3|C|+|I|=12$ slots. The element
\[
    \omega=6
\]
has order $12$ in $\F_{97}^\times$. Hence
\[
    \Omega=\set{1,\omega,\omega^2,\ldots,\omega^{11}}
\]
with actual values
\[
    \Omega=\set{1,6,36,22,35,16,96,91,61,75,62,81}.
\]
The vanishing polynomial is
\[
    Z_\Omega(X)=X^{12}-1.
\]

\subsection{Circuit and statement}

Take the three-gate computation
\[
    (x_1+x_2)(x_2+w_1).
\]
Let
\[
    x_1=5,
    \qquad
    x_2=6,
    \qquad
    w_1=1.
\]
Then
\[
    (5+6)(6+1)=11\cdot7=77.
\]
The computation trace is:
\[
\begin{array}{c|ccc|c}
\toprule
\text{Gate} & L & R & O & \text{Type}\\
\midrule
0 & 5 & 6 & 11 & +\\
1 & 6 & 1 & 7 & +\\
2 & 11 & 7 & 77 & \times\\
\bottomrule
\end{array}
\]

\subsection{Trace polynomial}

We place public and witness inputs at negative powers and gate wires at consecutive triples:
\[
\begin{array}{c|c|c|c}
\toprule
\text{Exponent }i & \omega^i & \text{Meaning} & T(\omega^i)\\
\midrule
11 & 81 & x_1 & 5\\
10 & 62 & x_2 & 6\\
9 & 75 & w_1 & 1\\
0 & 1 & \text{gate 0 left} & 5\\
1 & 6 & \text{gate 0 right} & 6\\
2 & 36 & \text{gate 0 output} & 11\\
3 & 22 & \text{gate 1 left} & 6\\
4 & 35 & \text{gate 1 right} & 1\\
5 & 16 & \text{gate 1 output} & 7\\
6 & 96 & \text{gate 2 left} & 11\\
7 & 91 & \text{gate 2 right} & 7\\
8 & 61 & \text{gate 2 output} & 77\\
\bottomrule
\end{array}
\]

There is a unique polynomial of degree at most $11$ interpolating these values. It is
\[
\begin{aligned}
T(X)={}&20+85X+22X^2+20X^3+80X^4+10X^5\\
&+47X^6+38X^7+27X^8+70X^9+6X^{10}+65X^{11}
\pmod{97}.
\end{aligned}
\]

\subsection{Public input check}

The verifier knows $x_1=5$ and $x_2=6$. Define $v(X)$ by
\[
    v(\omega^{11})=5,
    \qquad
    v(\omega^{10})=6.
\]
The interpolating degree-$1$ polynomial is
\[
    v(X)=45+51X\pmod{97}.
\]
The public input condition is
\[
    T(y)-v(y)=0
    \qquad
    \forall y\in\Omega_{\mathrm{inp}}=\set{\omega^{11},\omega^{10}}.
\]
Indeed,
\[
    T(81)=5=v(81),
    \qquad
    T(62)=6=v(62).
\]
This would be proved by ZeroTest on the two-point input domain.

\subsection{Gate check}

The gate-start domain is
\[
    \Omega_{\mathrm{gates}}=\set{\omega^0,\omega^3,\omega^6}=\set{1,22,96}.
\]
Define a selector polynomial $S$ by
\[
    S(1)=1,
    \qquad
    S(22)=1,
    \qquad
    S(96)=0.
\]
The degree-$2$ interpolant is
\[
    S(X)=68+49X+78X^2\pmod{97}.
\]
The gate polynomial is
\[
    G_{\mathrm{gate}}(X)
    =S(X)(T(X)+T(\omega X))+(1-S(X))T(X)T(\omega X)-T(\omega^2X).
\]
We need
\[
    G_{\mathrm{gate}}(y)=0
    \qquad\forall y\in\Omega_{\mathrm{gates}}.
\]
Concretely:
\[
\begin{array}{c|c}
\toprule
\text{Gate-start }y & \text{Check}\\
\midrule
1 & 5+6-11=0\\
22=\omega^3 & 6+1-7=0\\
96=\omega^6 & 11\cdot7-77=0\\
\bottomrule
\end{array}
\]
This condition is also proved by a ZeroTest, now over $\Omega_{\mathrm{gates}}$.

\subsection{Wiring as a prescribed permutation}

The copy constraints are:
\[
T(\omega^{10})=T(\omega^1)=T(\omega^3)=6,
\]
\[
T(\omega^{11})=T(\omega^0)=5,
\]
\[
T(\omega^2)=T(\omega^6)=11,
\]
\[
T(\omega^9)=T(\omega^4)=1.
\]
The permutation $W:\Omega\to\Omega$ acts by cycles
\[
    (\omega^{10}\ \omega^1\ \omega^3)
    (\omega^{11}\ \omega^0)
    (\omega^2\ \omega^6)
    (\omega^9\ \omega^4),
\]
with $\omega^5,\omega^7,\omega^8$ fixed. The logical wiring condition is
\[
    T(y)=T(W(y))
    \qquad\forall y\in\Omega.
\]

Interpolating $W$ on $\Omega$ gives one polynomial representative
\[
\begin{aligned}
W(X)={}&53X+38X^2+50X^3+8X^4+81X^5+43X^6\\
&+80X^7+94X^8+44X^9+21X^{10}+54X^{11}
\pmod{97}.
\end{aligned}
\]
A direct check of $T(X)-T(W(X))$ would involve a high-degree composition: the polynomial $W(X)$ is merely an interpolation of the wiring table, not a cheap function to compose into $T$. Plonk therefore uses a prescribed permutation ProductCheck.

\subsection{ProductCheck inside the wiring argument}

The prescribed permutation check samples random challenges $\rho,\sigma\in\F_{97}$ and proves
\[
    \prod_{a\in\Omega}
    \frac{\rho-\sigma W(a)-T(a)}{\rho-\sigma a-T(a)}=1.
\]
For this concrete example choose
\[
    \rho=17,
    \qquad
    \sigma=3.
\]
Define
\[
    F(a)=\rho-\sigma W(a)-T(a),
    \qquad
    G(a)=\rho-\sigma a-T(a),
\]
and
\[
    h(a)=\frac{F(a)}{G(a)}.
\]
The following table lists the domain order and the running product values
\[
    t(\omega^s)=\prod_{i=0}^{s}h(\omega^i).
\]

\begin{center}
\small
\begin{tabular}{c|c|c|c|c|c}
\toprule
$i$ & $a=\omega^i$ & $T(a)$ & $W(a)$ & $h(a)$ & $t(a)$\\
\midrule
0 & 1 & 5 & 81 & 39 & 39\\
1 & 6 & 6 & 22 & 91 & 57\\
2 & 36 & 11 & 96 & 37 & 72\\
3 & 22 & 6 & 62 & 12 & 88\\
4 & 35 & 1 & 75 & 83 & 29\\
5 & 16 & 7 & 16 & 1 & 29\\
6 & 96 & 11 & 36 & 21 & 27\\
7 & 91 & 7 & 91 & 1 & 27\\
8 & 61 & 77 & 61 & 1 & 27\\
9 & 75 & 1 & 35 & 90 & 5\\
10 & 62 & 6 & 6 & 66 & 39\\
11 & 81 & 5 & 1 & 5 & 1\\
\bottomrule
\end{tabular}
\end{center}

The final value is
\[
    t(\omega^{11})=1,
\]
so the product identity holds. Notice that the entries with $h(a)=1$ correspond to fixed points of the wiring permutation. The interpolated running product polynomial is
\[
\begin{aligned}
t(X)={}&69+25X+27X^2+18X^3+79X^4+64X^5\\
&+83X^6+17X^7+13X^8+88X^9+53X^{10}+85X^{11}
\pmod{97}.
\end{aligned}
\]

Now the prover interpolates the running product polynomial $t(X)$ from the values above. The recurrence to be checked is the rational ProductCheck recurrence:
\[
    t(\omega X)G(\omega X)=t(X)F(\omega X)
    \qquad\forall X\in\Omega.
\]
Equivalently, define
\[
    R(X)=t(\omega X)G(\omega X)-t(X)F(\omega X).
\]
Then
\[
    R(X)=0\qquad\forall X\in\Omega,
\]
so
\[
    q(X)=\frac{R(X)}{X^{12}-1}
\]
is a polynomial. A verifier checks this quotient relation at one random point $z$.

For example, take
\[
    z=2.
\]
The interpolated polynomials give
\[
    R(2)=25,
    \qquad
    z^{12}-1=2^{12}-1\equiv21\pmod{97}.
\]
Since
\[
    21^{-1}\equiv 37\pmod{97},
\]
we get
\[
    q(2)=R(2)(z^{12}-1)^{-1}=25\cdot37\equiv52\pmod{97}.
\]
Thus the verifier's random-point check is
\[
    R(2)
    \stackrel{?}{=}
    q(2)(2^{12}-1),
\]
and indeed
\[
    52\cdot21=1092\equiv25\pmod{97}.
\]
This is the ProductCheck embedded inside the Plonk wiring proof. The table proves the randomized pair-multiset equality. By the prescribed permutation lemma, this implies $T(a)=T(W(a))$ for every wire slot $a\in\Omega$, except with the usual random-challenge soundness error.

\subsection{Output check}

In this toy trace, the final output slot is
\[
    T(\omega^8)=77.
\]
If the statement is ``the circuit output is $77$,'' the verifier asks for a direct opening proving
\[
    T(\omega^8)=77.
\]
If the proof system uses a zero-output convention, one appends a final linear gate computing
\[
    T(\omega^8)-77
\]
and checks that the new final output is zero.

\subsection{The final proof obligations}

The complete Plonk proof for this toy circuit consists of the following obligations:
\begin{enumerate}[leftmargin=2em]
    \item \textbf{Trace commitment:} the prover commits to $T$.
    \item \textbf{Public inputs:} $T-v$ vanishes on $\Omega_{\mathrm{inp}}$.
    \item \textbf{Gate semantics:} $G_{\mathrm{gate}}$ vanishes on $\Omega_{\mathrm{gates}}$.
    \item \textbf{Wiring:} the prescribed permutation relation is proved by the rational ProductCheck above.
    \item \textbf{Output:} the final output equals $77$ or equals zero after a final adjustment gate.
\end{enumerate}

Each obligation is a low-degree polynomial identity checked at one or a few random points. A PCS such as KZG or Bulletproofs turns these oracle checks into a succinct argument \cite{KateZaveruchaGoldberg2010,Bulletproofs2018}. This is the complete architecture of Plonk in miniature.


